\label{ch:inference-process}

This chapter documents the statistical inference machinery: the 3-flavour
vacuum oscillation probability every model below shares
(\cref{sec:oscillation-probability}), the three composed model variants
(\cref{sec:model-variants}), the likelihood statistics
(\cref{sec:likelihoods}), the gradient-based fit procedure and its
uncertainty estimate (\cref{sec:fit-procedure}), and the 2-D likelihood
scan (\cref{sec:likelihood-scan}).

\section{Real 3-flavour vacuum oscillation probability}
\label{sec:oscillation-probability}

\modulehead{medium/vacuum/oscillation\_model.py}{Shared, detector-agnostic
3-flavour vacuum oscillation model (\pyclass{VacuumOscillationModel}),
reused unchanged from the rest of \texttt{tpeanuts}.}

\begin{theorybox}
Reactor antineutrinos at Daya Bay's real baselines ($\lesssim2$~km) cross
negligible Earth matter, so the exact coherent 3-flavour \emph{vacuum}
survival probability is used,
\begin{equation}
  P_{\bar\nu_e\bar\nu_e}(E,L;\theta) = \left|\sum_i U_{ei}^{*}\,
    e^{-i\Delta m^2_{i1}L/2E}\,U_{ei}\right|^2,
  \label{eq:vacuum-survival}
\end{equation}
with $U$ the PMNS matrix \parencite{Pontecorvo1957,MakiNakagawaSakata1962}
and antineutrino sign convention ($U\to U^*$ relative to the neutrino
case). Free parameters throughout this document: $\theta_{13}$ and $\Delta
m^2_{3l}\equiv\Delta m^2_{31}$ (normal ordering); fixed at Daya Bay's own
real external solar input, $\theta_{12}=\tfrac12\arcsin\sqrt{\sin^2
2\theta_{12}}$, $\Delta m^2_{21}$ (\code{survival\_probability\_solar.yaml},
\cref{ch:data-release}). $\theta_{23}$/$\delta_{CP}$ do not affect
\labelcref{eq:vacuum-survival} (a direct consequence of the standard
3-flavour PMNS parametrisation) and are held at an arbitrary NuFIT preset
value \parencite{EstebanNuFit2024}.
\end{theorybox}

\begin{implbox}
\pyclass{VacuumOscillationModel.predict\_pee(theta, L\_km, E\_grid\_MeV)}
implements \labelcref{eq:vacuum-survival} exactly (full 3-flavour, not a
2-flavour approximation), called once per (reactor, detector) pair by
\pyclass{DayaBayDetectorModel.predict} (\cref{sec:model-variants}).
\end{implbox}

\section{The three model variants}
\label{sec:model-variants}

\modulehead{detector/dayabay/inference\_model.py}{\pyclass{DayaBayDetectorModel},
\pyclass{JointDayaBayModel}, \pyclass{NearFarRatioDayaBayModel},
\pyclass{DayaBayExperimentalModel}: composition layer wrapping
\pyclass{VacuumOscillationModel} around \texttt{event\_rate.ibd\_event\_rate}
(\cref{ch:detector-implementation}).}

\subsection*{Per-detector and joint absolute-count models}

\begin{theorybox}
\pyclass{DayaBayDetectorModel} evaluates \labelcref{eq:vacuum-survival}
once per real reactor (6 calls, each at that reactor's own real baseline),
then calls \labelcref{eq:dayabay-full-prediction} to fold the real
6-reactor flux sum through the real cross section, response, and
background. Its free-parameter vector is assembled, in order, from
\begin{equation}
  \theta = \bigl(\underbrace{\theta_{13},\,\Delta m^2_{3l}}_{\rm oscillation}\bigr)
  \oplus \bigl[\nu_{\rm norm}\bigr]
  \oplus \bigl[c_1,\dots,c_5\bigr]
  \oplus \bigl[a_1,\dots,a_4\bigr],
  \label{eq:theta-assembly}
\end{equation}
with each bracketed group present only if the corresponding
\code{normalization\_free}/\code{background\_free}/\code{lsnl\_free} flag
is set. \pyclass{JointDayaBayModel} concatenates several
\pyclass{DayaBayDetectorModel} predictions (typically all 8 real
detectors) sharing one $\theta$, for a full combined fit -- the
\textbf{full-spectral-shape} strategy.
\end{theorybox}

\subsection*{Rate-only model}

\begin{theorybox}
A rate-only observable integrates \labelcref{eq:dayabay-full-prediction}
over the full real analysis-bin range, discarding the intra-detector
spectral shape but retaining the real detector-to-detector rate deficit:
\begin{equation}
  \keyresult{N_d^{\rm rate}(\theta) = \sum_k N_k^{(d)}(\theta),}
  \label{eq:rate-only-observable}
\end{equation}
one real number per detector. Because a rate-only observable has
essentially no independent $\Delta m^2_{31}$ leverage in this two-baseline-tier
geometry, $\Delta m^2_{31}$ is fixed at Daya Bay's own published value and
only $\theta_{13}$ (plus $\nu_{\rm norm}$) is fit.
\end{theorybox}

\subsection*{Near/far ratio model}

\begin{theorybox}
Mirroring Daya Bay's own original 2012 discovery strategy
\parencite{An2012DayaBayDiscovery}, the far-hall/near-hall per-bin count
ratio is predicted instead of absolute counts:
\begin{equation}
  \keyresult{R_k(\theta) = \frac{N_k^{\rm far}(\theta)}{N_k^{\rm near}(\theta)},
    \qquad N_k^{\rm near,far}(\theta) = \!\!\sum_{d\,\in\,{\rm NEAR,FAR}}\!\! N_k^{(d)}(\theta),}
  \label{eq:nearfar-ratio}
\end{equation}
with $\mathrm{NEAR}=\{\mathrm{AD11,AD12,AD21,AD22}\}$ (EH1+EH2) and
$\mathrm{FAR}=\{\mathrm{AD31,\dots,AD34}\}$ (EH3,
\cref{sec:detector-geometry}). This cancels, by construction, every
systematic common to all 8 real detectors: with \code{normalization\_free
=False}, $\nu_{\rm norm}$ is not even a free parameter of this model (it
would cancel in \labelcref{eq:nearfar-ratio} to the extent the real
background is negligible relative to signal, and only approximately once
it is not); any \emph{residual} mismatch in the simplified IAV/LSNL/cross-section
response also cancels approximately, since it multiplies the near and far
prediction identically.
\end{theorybox}

\subsection*{Experimental reparametrisation}

\begin{theorybox}
Daya Bay's own combined analyses report $\sin^2(2\theta_{13})$ and the
effective splitting
\begin{equation}
  \Delta m^2_{ee} = \cos^2\theta_{12}\,\Delta m^2_{31} + \sin^2\theta_{12}\,\Delta m^2_{32}
  \label{eq:dm2ee-definition}
\end{equation}
\parencite{NunokawaParkeZukanovich2005,An2023DayaBayPrecision}, not the bare
PMNS $(\theta_{13},\Delta m^2_{3l})$ this project's models natively expose.
With $\theta_{12}$/$\Delta m^2_{21}$ fixed, substituting $\Delta
m^2_{32}=\Delta m^2_{31}-\Delta m^2_{21}$ into
\labelcref{eq:dm2ee-definition} and solving gives an exact, fixed linear
offset,
\begin{equation}
  \keyresult{\Delta m^2_{31} = \Delta m^2_{ee} + \sin^2\theta_{12}\,\Delta m^2_{21},}
  \label{eq:dm31-from-dm2ee}
\end{equation}
so \pyclass{DayaBayExperimentalModel} is a pure, exactly invertible change
of variables (verified in this project's own test suite), not a different
fit -- it exists purely so a fit's free-parameter vector and any resulting
contour are directly comparable to Daya Bay's own published numbers/plots.
\end{theorybox}

\begin{implbox}
\begin{itemize}
  \item \pyclass{DayaBayDetectorModel(oscillation\_model, detector,
    normalization\_free=True, background\_free=False, lsnl\_free=False)}:
    \code{.free} implements \labelcref{eq:theta-assembly};
    \code{.predict(theta)} slices $\theta$ in that same order and calls
    \pyfunc{ibd\_event\_rate}.
  \item \pyclass{JointDayaBayModel(models)}: \code{.predict} concatenates.
  \item \pyclass{NearFarRatioDayaBayModel(near\_models, far\_models)}
    implements \labelcref{eq:nearfar-ratio};
    \code{.real\_observed\_ratio()} (a \code{@staticmethod}) returns the
    real observed ratio and its propagated uncertainty
    (\cref{sec:likelihoods}).
  \item \pyclass{DayaBayExperimentalModel(model, theta12, dm21)} implements
    \labelcref{eq:dm31-from-dm2ee}, converting \code{theta}'s first two
    entries before delegating to the wrapped model.
\end{itemize}
\end{implbox}

\section{Likelihood statistics}
\label{sec:likelihoods}

\modulehead{inference/likelihood.py}{\pyfunc{poisson\_nll},
\pyfunc{chi2\_asymmetric}, \pyfunc{gaussian\_prior\_penalty} -- every
$-2\ln L$ statistic used in this document.}

\begin{theorybox}
\subsection*{Poisson (Baker-Cousins) likelihood}
Absolute per-bin counts (full-shape and rate-only fits) are compared via
the Poisson likelihood-ratio statistic \parencite{BakerCousins1984},
\begin{equation}
  \keyresult{-2\ln L_{\rm Poisson} = 2\sum_k\Bigl[N_k-n_k^{\rm obs}+n_k^{\rm obs}
    \ln\frac{n_k^{\rm obs}}{N_k}\Bigr],}
  \label{eq:poisson-nll}
\end{equation}
with $N_k$/$n_k^{\rm obs}$ the predicted/observed counts. The extra
$-n_k^{\rm obs}+n_k^{\rm obs}\ln n_k^{\rm obs}$ terms subtract the
saturated model's own likelihood, so \labelcref{eq:poisson-nll} is exactly
0 when $N_k=n_k^{\rm obs}$ bin by bin and behaves like an ordinary
chi-square asymptotically (Wilks' theorem).

\subsection*{Gaussian (chi2\_asymmetric) likelihood}
The near/far ratio (a ratio of counts, not itself Poisson-distributed) is
compared via a Gaussian statistic with one-sided uncertainty,
\begin{equation}
  -2\ln L_{\chi^2} = \sum_k \left(\frac{R_k-R_k^{\rm obs}}{\sigma_k}\right)^2,
  \label{eq:chi2-asymmetric}
\end{equation}
with the real observed ratio's uncertainty propagated from independent
Poisson numerator/denominator counts,
\begin{equation}
  \keyresult{\sigma_k = R_k^{\rm obs}\sqrt{\frac{1}{N_k^{\rm far,obs}}+\frac{1}{N_k^{\rm near,obs}}}.}
  \label{eq:ratio-sigma-propagation}
\end{equation}

\subsection*{Gaussian prior penalty}
Real, externally measured nuisance-parameter constraints (background-rate
scales, LSNL pulls, \cref{sec:nuisances}) are added as an independent
Gaussian penalty on the free-parameter vector itself, not the prediction:
\begin{equation}
  \keyresult{-2\ln L_{\rm penalty} = \sum_j \left(\frac{\theta_j-\mu_j}{\sigma_j}\right)^2,}
  \label{eq:gaussian-prior-penalty}
\end{equation}
added to \labelcref{eq:poisson-nll} (or \labelcref{eq:chi2-asymmetric}) to
give the total statistic minimized.
\end{theorybox}

\begin{implbox}
\begin{itemize}
  \item \pyfunc{poisson\_nll(prediction, value)} implements
    \labelcref{eq:poisson-nll}, built as
    \code{xlogy(value,value)-xlogy(value,prediction)} rather than the
    algebraically equal single-\code{xlogy} form, to keep the backward
    pass finite at an empty bin (\code{value==0}).
  \item \pyfunc{chi2\_asymmetric(prediction, value, sigma\_minus, sigma\_plus)}
    implements \labelcref{eq:chi2-asymmetric}.
  \item \pyfunc{gaussian\_prior\_penalty(values, prior\_mean, prior\_sigma)}
    implements \labelcref{eq:gaussian-prior-penalty}; passed as
    \pyfunc{fit\_lbfgs}'s \code{penalty\_fn} argument
    (\cref{sec:fit-procedure}).
  \item \pyfunc{NearFarRatioDayaBayModel.real\_observed\_ratio} implements
    \labelcref{eq:ratio-sigma-propagation}.
\end{itemize}
\end{implbox}

\section{Gradient-based fit procedure}
\label{sec:fit-procedure}

\modulehead{inference/fit.py}{\pyfunc{fit\_lbfgs}/\pyfunc{minimize\_lbfgs}:
LBFGS point estimate plus Laplace/Fisher covariance, for any
\pyclass{DifferentiableModel}.}

\begin{theorybox}
\subsection*{Parameter rescaling}
LBFGS's identity-initialised inverse-Hessian approximation assumes roughly
comparable parameter scales; $\theta_{13}\sim\mathcal O(1)$~rad and
$\Delta m^2_{31}\sim\mathcal O(10^{-3})$~eV$^2$ differ by 3-4 orders of
magnitude, which starves the line search of a usable step size. A
dimensionless $u=\theta/{\rm scale}$ is fit instead, with ${\rm scale}$ set
from the loss gradient at $\theta_0$,
\begin{equation}
  {\rm scale} = \bigl|\nabla_\theta(-2\ln L)\bigr|_{\theta_0}^{-1},
\end{equation}
undone only in the returned $\hat\theta$ -- every model always sees
physical-unit $\theta$.

\subsection*{Optimisation and convergence}
$\hat\theta=\arg\min_\theta\bigl[-2\ln L(\theta) + {\rm penalty}(\theta)\bigr]$
is found by \code{torch.optim.LBFGS} using the exact autograd gradient of
the total statistic through every real physics step of
\cref{ch:detector-implementation} (no finite differences).

\subsection*{Laplace/Fisher covariance}
The marginal uncertainty on each fitted parameter is the Laplace
approximation: with $\mathcal I=\tfrac12\nabla^2_\theta(-2\ln L)$ the
Fisher information at the optimum (the autograd Hessian of the total,
penalised statistic),
\begin{equation}
  \keyresult{\Sigma = \mathcal I^{-1}, \qquad \sigma_j = \sqrt{\Sigma_{jj}}.}
  \label{eq:laplace-covariance}
\end{equation}
A direction the data barely constrains gives a Hessian eigenvalue near
zero (occasionally slightly negative from numerical round-off); small or
negative eigenvalues are floored at a tiny fraction of the largest one
before inversion, so an unconstrained direction gets a large but finite
$\sigma_j$ instead of \code{nan}.
\end{theorybox}

\begin{implbox}
\begin{itemize}
  \item \pyfunc{fit\_lbfgs(model, theta0, value, sigma\_minus=None,
    sigma\_plus=None, likelihood=..., penalty\_fn=None, max\_iter=200)}
    returns a \pyclass{FitResult} (\code{theta\_hat}, \code{param\_names},
    \code{chi2\_history}, \code{covariance}, and a \code{sigma} property).
  \item \code{penalty\_fn}: an optional \code{Callable[[theta],\ Tensor]}
    added to the loss inside the LBFGS closure \emph{and} to the Hessian
    computation for \labelcref{eq:laplace-covariance} -- a tight prior
    correctly shrinks its own parameter's reported uncertainty. \code{None}
    (default) reproduces the previous, unconstrained-fit behaviour exactly.
\end{itemize}
\end{implbox}

\section{2-D likelihood scan}
\label{sec:likelihood-scan}

\modulehead{inference/scan.py}{\pyfunc{loglik\_grid}: evaluates $-2\ln L$
on a $(\theta_x,\theta_y)$ grid, for visual/contour diagnostics.}

\begin{theorybox}
On a chosen 2-D grid, every entry of \code{model.free} other than the two
scanned parameters is held fixed at a baseline $\theta$ (a cheap slice
through the likelihood surface) or re-minimised at each grid point (the
statistically correct construction when other parameters are also
uncertain, at the cost of one LBFGS run per grid point). With exactly two
free parameters in \code{model.free} (the rate-only and near/far models,
once their own nuisances are fixed/absent), nothing remains to profile and
both modes coincide -- the scans in \cref{ch:results} are genuine 2-D
likelihood surfaces of the model's full free set, not fixed-nuisance
slices.
\end{theorybox}

\begin{implbox}
\pyfunc{loglik\_grid(model, theta\_baseline, param\_x, x\_grid, param\_y,
y\_grid, value, sigma\_minus=None, sigma\_plus=None, likelihood=...,
profile\_others=False)} returns a real tensor shaped \code{(nx,ny)};
subtracting its minimum and comparing against standard 2-D $\Delta\chi^2$
confidence levels draws confidence contours.
\end{implbox}
